% !TEX root = ../main.tex
\chapter{Banks and Global Liquidity}\label{ch:globalliq}
\begin{paracol}{2}

\begin{sources}
\begin{tabularx}{\linewidth}{@{}l l L@{}}
\toprule
\textbf{Segment} & \textbf{Length} & \textbf{Content}\\
\midrule
\seg{15}{1}{L15.1 FT: European money market funds} & 2:37 & EUR~100 billion moved across currency areas.\\
\seg{15}{2}{L15.2 International transactions} & 11:42 & Bagehot's world, international: bills discounted in London, two FX dealers, the Bank of England.\\
\seg{15}{3}{L15.3 Dealer model for foreign exchange} & 10:05 & Gold points as the outside spread of FX dealers.\\
\seg{15}{4}{L15.4 Central banking, defense of domestic exchange} & 8:43 & The deficit country's central bank as dealer: pay out, sell assets, borrow reserves.\\
\seg{15}{5}{L15.5 Bank of England, defense against external drain} & 12:08 & Bank rate as outside spread; external drain; raising Bank rate; suspension.\\
\seg{15}{6}{L15.6 Toward a theory of exchange} & 9:37 & Without gold both outside spreads vanish; Keynes's \emph{Tract} (1923) and covered interest parity.\\
\bottomrule
\end{tabularx}

\smallskip
\textbf{Files.} Transcripts \file{materials/transcripts/L15-P*.txt}; Mehrling's notes \mn{15}{1} (headed ``18.'' in the PDF, from an earlier numbering of the course).
\end{sources}

\section*{Motivation}
Today's world is one of flexible, volatile exchange rates (next lecture). This lecture goes back a hundred years, to the gold standard, to introduce concepts
needed next time \ts{15}{1}{2:10}. It is ``the world that Bagehot knew'' (lecture~\ref{ch:bagehot}) again, now with its international piece: firms in different
countries trading by bills of exchange discounted in London, foreign exchange dealers, and the Bank of England on top \ts{15}{2}{0:03}. The tool is
Treynor's dealer model (lecture~\ref{ch:dealers}), applied twice: to exchange rates and to interest rates.

% ------------------------------------------------------------------
\section{FT: European money market funds}\label{sec:gl-ft}

\begin{casestudy}[FT, 2012: ``Europe's money market funds look further afield'']
Citing a Fitch Ratings report: rating downgrades, worries over the eurozone and lower bank demand for short-term funding are pushing European money market
funds towards Asia and the Middle East. Over two years they moved almost a fifth of their assets, about EUR~100 billion, from the UK, the US and the eurozone
periphery to Germany, France, the Netherlands, the Nordics, Asia and the Middle East \ts{15}{1}{0:48}. Moving EUR~100 billion
across currency areas has exchange-rate consequences: part of the volatility of flexible rates comes from such flows, and from people trying to avoid
currencies that are going to depreciate \ts{15}{1}{1:29}.
\end{casestudy}

\begin{aside}[Coming attractions]
The same day's FT reported China lifting foreign access to its markets, linked to the week's reading by McCauley (BIS) on the internationalization of the renminbi,
discussed next lecture \ts{15}{1}{0:07}.
\end{aside}

% ------------------------------------------------------------------
\section{International transactions in Bagehot's world}\label{sec:gl-transactions}

\paragraph{The bill discounted in London.} A surplus firm (or country) sells goods to a deficit firm and receives a bill of exchange, a promise to pay in 90 days;
neither is in England \ts{15}{2}{0:44}. The bill is not money, so the seller takes it to a private bank in the City of London,
the centre of the world money market in 1873, which discounts it, accepting the bill and paying out sterling notes \ts{15}{2}{1:50}.
(Notes rather than deposits, because the intuition is sharper \ts{15}{2}{2:54}.) The surplus firm is paid before the deficit firm has paid. In 90 days the
deficit firm pays the bill with notes, which flow back to the bank \ts{15}{2}{3:14}. The banks manage the outflow and inflow of
notes, holding bills that mature every day: classic 19th-century bill banking \ts{15}{2}{3:56}.

\paragraph{The foreign exchange dimension.} Sterling notes are high-powered money for international banking, but they cannot pay the surplus firm's workers;
and the deficit firm sells its goods for its own currency and must turn it into sterling to pay the bill \ts{15}{2}{4:38}. So
there are two foreign exchange dealers (or banks: in this course a bank is a special kind of dealer and a dealer a special kind of bank)
\ts{15}{2}{5:39}:
\begin{itemize}
  \item the \textbf{surplus dealer} buys sterling notes from the surplus firm and pays out domestic currency \ts{15}{2}{6:01};
  \item the \textbf{deficit dealer}, 90 days later, sells sterling notes to the deficit firm for its domestic currency \ts{15}{2}{7:06}.
\end{itemize}
The two domestic currencies may differ; what the dealers have in common is that both exchange against the international currency, sterling \ts{15}{2}{8:10}.

\paragraph{The Bank of England.} In 1873 sterling notes are liabilities of the Issue Department of the Bank of England, which holds gold and keeps the notes at par
with it \ts{15}{2}{8:31}. Both dealers use Bank of England notes as international reserves, and may cash them for gold or vice versa: the Bank makes a two-way
market between notes and gold \ts{15}{2}{8:52}. Three kinds of dealers, in a hierarchy: gold at the top, Bank of
England notes promising gold, then the other currencies kept at par with those notes \ts{15}{2}{9:58}.

\begin{nfigure}
\begin{taccts}
\tacct[2.2cm]{Surplus firm}{$-$Goods\\$+$Bill, then\\$-$Bill $+$Notes}{}\quad
\tacct[2.2cm]{City bank}{$+$Bill\\(90 days later: $-$Bill)}{$-$Notes (paid out)\\(90 days later: $+$Notes)}\quad
\tacct[2.2cm]{Deficit firm}{$+$Goods}{$+$Bill (due in 90 days)}
\end{taccts}
\begin{taccts}
\tacct[2.2cm]{Surplus dealer}{$+$Sterling notes}{$+$Domestic currency}\quad
\tacct[2.2cm]{Deficit dealer}{$-$Sterling notes}{$-$Domestic currency}\quad
\tacct[2.2cm]{Bank of England (Issue Dept.)}{$\pm$Gold}{$\pm$Notes}
\end{taccts}
\caption{International trade finance in 1873, as drawn in class \ts{15}{2}{1:08}.}\label{fig:gl-taccts}
\end{nfigure}

This is how it worked, and how it works today with the dollar: most world trade is invoiced in dollars, though much of it does not involve the US; back then it
was sterling \ts{15}{2}{10:59}.

% ------------------------------------------------------------------
\section{The dealer model for foreign exchange}\label{sec:gl-fxdealer}

Recall from lecture~\ref{ch:chartalism}: if the dollar is $x$ ounces of gold and the pound $y$ ounces, the mint parity exchange rate, quoted in pounds per dollar
(the convention of this course), is $x/y$ \ts{15}{3}{0:44}. The actual rate can deviate, because moving gold is costly:
\begin{keyformulas}
\textbf{Gold points} \ts{15}{3}{1:46}: with $\delta$ set by the cost of shipping gold,
\[
  \frac{x}{y}-\delta \;\le\; E \;\le\; \frac{x}{y}+\delta \qquad (E \text{ in pounds per dollar}).
\]
\end{keyformulas}

Inside the band there is no arbitrage to snap the rate back. This is the Treynor model hinted at two lectures ago: the gold points are an \emph{outside spread},
and the FX dealers do business inside it \ts{15}{3}{2:26}.

\paragraph{The deficit country's dealer.} Let the deficit country be the US (deliberately: the world currency is the pound, not the dollar; ``God'' did not plan
the dollar to be the reserve currency, certainly not in 1873) \ts{15}{3}{0:24}. Its dealers buy dollars and sell sterling: long the
domestic currency, short the international one, losing reserves, so long liquidity risk \ts{15}{3}{3:50}. They do it only at a lower price of the
dollar: with pounds per dollar, a lower number is a depreciated dollar \ts{15}{3}{5:13}. They take exchange risk only if paid in
expectation, hoping the dollar snaps back so they can sell high \ts{15}{3}{6:14}.

\paragraph{Why the dealer dares.} What if the dollar keeps falling? It cannot fall far: dollars can be taken to the central bank and converted into gold, which can
be shipped to England for sterling notes. The gold point is a floor, the outside spread; the dealers' quotes are the inside spread
\ts{15}{3}{6:55}. The quotes seen in the market are the dealers'; everyone knows the mint
pars, and a large volume of business pushes the price to the gold points, where gold flows keep it within bounds as long as central banks can supply the
gold \ts{15}{3}{8:37}.

\begin{coursenote}
The US central bank in this example is hypothetical: in 1873 the US had no central bank (the Fed dates from 1913), as Mehrling remarks \ts{15}{3}{7:36}.
\end{coursenote}

\begin{nfigure}
\begin{tikzpicture}[font=\scriptsize,x=1cm,y=1cm]
  % left panel: FX dealer
  \begin{scope}
  \draw[-{Stealth[length=4pt]}] (-3.2,0) -- (3.4,0) node[below left,align=right]{dealer's dollar\\inventory};
  \draw[-{Stealth[length=4pt]}] (0,-0.2) -- (0,4.4) node[above,align=center]{$E$ (\pounds/\$)};
  \draw[very thick,thmcol] (-3.2,3.4) -- (3.2,3.4) node[above left]{$x/y+\delta$ (gold point)};
  \draw[very thick,thmcol] (-3.2,0.8) -- (3.2,0.8) node[below left]{$x/y-\delta$ (gold point)};
  \draw[dotted] (-3.2,2.1) -- (3.2,2.1); \node[left] at (-3.2,2.1) {$x/y$};
  \draw[very thick,defcol] (-3,3.4) -- (3,1.1);
  \draw[very thick,intcol] (-3,3.1) -- (3,0.8);
  \fill (1.8,1.4) circle (1.5pt);
  \draw[thin] (1.8,1.4) -- (2.2,2.6) node[above]{deficit country};
  \node[align=center] at (0,-1.25) {FX dealer: exchange rate space};
  \end{scope}
  % right panel: bank as dealer in yield space
  \begin{scope}[xshift=7.3cm]
  \draw[-{Stealth[length=4pt]}] (-3.2,0) -- (3.4,0) node[below left,align=right]{bank's bill\\holdings};
  \draw[-{Stealth[length=4pt]}] (0,-0.2) -- (0,4.4) node[above]{discount rate};
  \draw[very thick,thmcol] (-3.2,3.4) -- (3.2,3.4) node[above left]{Bank rate};
  \draw[very thick,defcol] (-3,1.1) -- (3,3.4);
  \draw[very thick,intcol] (-3,0.8) -- (3,3.1);
  \node[align=center] at (1.6,1.3) {market rate};
  \node[align=center] at (0,-1.25) {City bank: interest rate (yield) space};
  \end{scope}
\end{tikzpicture}
\caption{Two outside spreads of the gold standard: the gold points for FX dealers \ts{15}{3}{3:28} and Bank rate for the discounting banks
\ts{15}{5}{1:26}. In yield space the curves slope up. Schematic.}\label{fig:gl-spreads}
\end{nfigure}

% ------------------------------------------------------------------
\section{The central bank defends the domestic exchange}\label{sec:gl-defense}

Suppose the dollar is pushed to the gold point: private dealers, who have as much risk as they want, bring domestic currency to the central bank and ask for
international currency, gold or sterling notes \ts{15}{4}{0:04}. Domestic currency is the central bank's own liability, so it is
simply cancelled: the central bank's balance sheet shrinks, it loses reserves and high-powered money contracts \ts{15}{4}{1:10}.
Textbooks say a currency is defended by contractionary monetary policy; here that follows literally from acting as a dealer, not from an intention
\ts{15}{4}{1:51}.

Its gold and notes are limited, so the central bank must think ahead. Three logical options, as in the flow-of-funds accounting
\ts{15}{4}{2:33}:
\begin{enumerate}
  \item \textbf{pay out reserves} (dishoarding);
  \item \textbf{sell an asset} for reserves, for example Treasury bills, which pushes bill prices down and interest rates up \ts{15}{4}{3:17};
  \item \textbf{borrow reserves} from another central bank, possibly the Bank of England itself \ts{15}{4}{3:58}.
\end{enumerate}

\begin{nfigure}
\begin{taccts}
\tacct[2.6cm]{Deficit country's central bank}{$-$Gold or sterling notes\\(1) $-$Reserves\\(2) $-$T-bills, $+$Reserves\\(3) $+$Reserves}{$-$Domestic currency\\\\(3) $+$Borrowed reserves}
\end{taccts}
\caption{Defending the domestic exchange at the gold point \ts{15}{4}{0:50}.}\label{fig:gl-cbdefense}
\end{nfigure}

The central bank is a backstop not so much because it holds large reserves, but because it can replenish them: it has assets (its special relationship with the
sovereign) and relationships with other central banks, which a private dealer lacks \ts{15}{4}{4:42}.

\begin{intuition}[What ``raising rates to defend the currency'' really means]
The notes read the three channels in balance sheets \mn{15}{3}\mn{15}{4}. Upward pressure on the currency is no problem: issue currency to buy reserves. Downward pressure means buying the
currency, the central bank's own liability, so its balance sheet contracts: paying out reserves works only while they last; selling Treasury bills pushes their price down and interest rates up,
perhaps a lot if speculators fear devaluation. ``We often hear loose talk about central banks defending their currency by raising interest rates'': what really happens is that the central bank
offers interest-bearing securities in exchange for currency, hoping the interest is enough to keep holders from demanding international reserves. Borrowing reserves from other central banks
(or the IMF, ``somewhat like an international central bank'') leaves the balance sheet's size unchanged: part of the monetary base in private hands shifts to other central banks. In every
case, contraction of the money supply supports the currency because the central bank buys its own liability. And the deficit entity that must adjust need not be below the Bank of England: it
can be the Bank itself \mn{15}{4}.
\end{intuition}

\begin{intuition}[What counts as reserves depends on the layer]
A student asks what the difference is between domestic currency and reserves. For the US central bank in a sterling world, reserves are gold or Bank of England
notes; for a bank inside the US, reserves would be domestic currency. What counts as reserves and what as credit depends on where you are in the hierarchy.
Foreign exchange is everything learned so far, one layer higher, with one difference: inside a country the relations are par relations, between countries the
rates move between the gold points \ts{15}{4}{5:44}.
\end{intuition}

\paragraph{The abstraction that removes liquidity.} Economists often treat the gold points as so narrow that the rate is just the mint par ratio, $\delta=0$, a
horizontal line. That abstracts from the essence of monetary economics: with no price movement no dealer would make a market, and without dealers there is no
liquidity. The course puts back the market mechanics that give liquidity \ts{15}{4}{7:28}.

% ------------------------------------------------------------------
\section{The Bank of England and the external drain}\label{sec:gl-boe}

\paragraph{City banks as dealers in yield space.} The City banks discount bills and wait 90 days for the money: they lend long and borrow short, take liquidity
risk, and so are dealers too, in interest-rate space; in yield space the bid and offer curves slope up, and the outside spread is \emph{Bank rate}
\ts{15}{5}{0:03}. A bank losing cash faster than it comes in raises its discount rate, so fewer
bills come in and the inflow exceeds the outflow \ts{15}{5}{1:47}. If all banks do it, the market rate of interest rises,
eventually to Bank rate, at which the Bank of England, as lender of last resort, rediscounts bills for notes \ts{15}{5}{2:28}.
The banks dare to discount because the backstop is there \ts{15}{5}{4:17}.

\paragraph{Internal drain.} The Banking Department pays out the notes it holds as its reserve, or it expands its balance sheet on both sides, crediting deposits:
as long as banks accept deposits as notes it economizes on its note issue and keeps rates from rising above Bank rate (lecture~\ref{ch:bagehot})
\ts{15}{5}{4:39}.

\paragraph{External drain.} New today: what if the dollar is bid \emph{up} to the other gold point? That outside spread is the Bank of England's market: surplus
dealers with too many notes ask for gold, and the Issue Department pays, $-$notes, $-$gold, defending the mint par of sterling
\ts{15}{5}{5:40}. Its gold is limited, so like any central bank it seeks more reserves. Its equivalent
of selling Treasury bills is to \textbf{raise Bank rate}, at a time when everyone is bringing bills for discount: the world rate of interest rises, borrowers
repay and do not renew, and gold flows in \ts{15}{5}{7:26}. But a high Bank rate can bankrupt people and cause
further problems \ts{15}{5}{8:50}.

\begin{definition}[New concepts: external drain, suspension of specie payments]
An \emph{internal drain} is a demand for the central bank's notes from the domestic system; an \emph{external drain} is a demand for gold to pay abroad
\ts{15}{5}{5:40}. \emph{Suspending specie payments}: the central bank stops converting its notes into gold (``specie'', coined metal, from Latin
\emph{in specie}, ``in kind'') while still issuing notes and deposits \ts{15}{5}{8:50}.
\end{definition}

\paragraph{The nuclear option.} The third option is suspension: ``I'm no longer going to pay notes in gold'', the system becomes a pure credit system for the time
being \ts{15}{5}{9:11}. Other central banks that suspend drop out of the international system; when the Bank of England does it, it
\emph{is} the international system: people keep using sterling as reserves, now unattached to gold, and perhaps decades later it returns to gold
\ts{15}{5}{9:51}.

\begin{keyformulas}
\textbf{The gold standard as discipline} \ts{15}{5}{10:53}. Domestically the central bank is at the apex and can
always print: no internal run can trouble it. Internationally an external drain forces it to defend its gold by raising Bank rate, at the cost of the domestic
economy, because foreigners did their banking business in London; the only alternative is suspension.
\end{keyformulas}

\begin{coursenote}
In class the comparison is said to be between ``lecture 10 and lecture 18''; it is between lecture~9 of these notes (``The World that Bagehot Knew'', which
Mehrling calls lecture~10) and this lecture~15 \ts{15}{5}{10:32}.
\end{coursenote}

\begin{history}[Suspensions of the Bank Charter Act]
Under the Bank Charter Act of 1844 the Issue Department's notes beyond a fixed fiduciary issue had to be backed by gold. In the crises of 1847, 1857 and 1866 the
government authorised the Bank to exceed that limit (a ``suspension'' of the Act, though convertibility was kept). Convertibility itself was suspended from 1797 to
1821 (the Restriction period) and, in effect, at the outbreak of the First World War; Britain returned to gold in 1925 and left it in 1931.\par
\emph{References:} J.~Clapham, \emph{The Bank of England: A History}, vol.~2, Cambridge UP, 1944; Bank of England, ``A brief history of banknotes''.
\end{history}

% ------------------------------------------------------------------
\section{Toward a theory of exchange}\label{sec:gl-theory}

\paragraph{After suspension.} A student: if payments are suspended, does the model not break down, and should sterling not be strongly devalued? Two answers
\ts{15}{6}{0:03}:
\begin{enumerate}
  \item yes, the story of mint par ratios and outside spreads goes away; it is replaced by something else, next lecture \ts{15}{6}{0:24};
  \item ``depreciate against what?'' Certainly against gold; but sterling may still be the best currency in the world, so not necessarily against other
        currencies: the whole currency system may depreciate against gold, but maybe not sterling against the rest \ts{15}{6}{1:05}.
\end{enumerate}

\paragraph{Two outside spreads gone.} The gold standard had two central bank outside spreads: the gold points, giving FX dealers room to work, and Bank rate,
giving discounting banks room to work \ts{15}{6}{2:08}. Moving away from gold removes both. Bank rate was a discount rate for 90-day bills; today
the Fed sets an overnight rate, the Fed funds rate, and does not set 90-day bill prices except in a crisis (when it walked up the yield curve)
\ts{15}{6}{2:49}. Next time the same two diagrams return, but the outside spreads will not be central banks
\ts{15}{6}{3:51}. Why explain the gold standard with the dealer model? Because, unlike the usual textbook explanation, it extends to fiat money: going off gold
does not mean ``la la land'' with no economics \ts{15}{6}{4:12}.

\paragraph{Keynes, 1923.} In \emph{A Tract on Monetary Reform}, his first real money book, written after the war, while countries were trying to get back to gold,
Keynes calls the gold standard ``already a barbarous relic'' (p.~172): central bankers now care about the stability of prices and employment and will not
sacrifice them to the ``outworn dogma'' of a gold price of \pounds3 17s 10\textonehalf d per ounce \ts{15}{6}{4:53}.
War finance had accustomed the world to paper money not backed by gold \ts{15}{6}{6:14}. His hope: each central bank stabilizes prices at home,
exchange rates float, and through purchasing power parity they are stable \ts{15}{6}{6:55}. It did not work: the world went back to gold, which
broke apart, and there was a world depression; this is the dream Mundell talked about \ts{15}{6}{7:17}.

But Keynes also started to discuss forward exchange, and wrote an equation equivalent to covered interest parity (lecture~\ref{ch:eurodollar}), the key
relationship for translating the two diagrams to a world without gold \ts{15}{6}{7:58}:
\begin{keyformulas}
\textbf{Covered interest parity} \ts{15}{6}{8:39}: with $S$ the spot and $F$ the forward rate (pounds per dollar), $R$ and $R^\ast$ the dollar and sterling term
rates of interest,
\[
  \frac{F}{S}=\frac{1+R^\ast}{1+R}.
\]
\end{keyformulas}

\begin{coursenote}
The auto-captions give the gold price as ``\pounds3 17 shillings 102''; the British mint price of standard gold was \pounds3 17s 10\textonehalf d per troy
ounce, as Keynes himself writes it, ``\pounds3:17:10\textonehalf\ per ounce''.

Checked against the scan (\file{materials/readings/L15-Keynes_Tract_on_Monetary_Reform.pdf}):
\begin{itemize}
  \item ``In truth, the gold standard is already a barbarous relic'' is on p.~172.
  \item ``Outworn dogma'' follows on p.~173, with the remark that a regulated non-metallic standard ``has slipped in unnoticed. \emph{It exists.}''
\end{itemize}
The equation is written on the board in class and is garbled in the captions; it is reproduced here in the notation of lecture~\ref{ch:eurodollar}.
\par\emph{Reference:} J.~M.~Keynes, \emph{A Tract on Monetary Reform}, Macmillan, 1923, ch.~3 (forward exchange) and ch.~4 (``barbarous relic'').
\end{coursenote}

% ------------------------------------------------------------------
\flushside
\section*{Key formulas and ideas}
\begin{keyformulas}
\begin{tabularx}{\linewidth}{@{}l L@{}}
Trade finance & bill discounted in London for sterling notes; FX dealers in each country exchange local currency for sterling\\
Hierarchy & gold $>$ Bank of England notes $>$ other currencies (kept within the gold points)\\
Gold points & $x/y-\delta \le E \le x/y+\delta$: outside spread for FX dealers\\
Defense & deficit central bank: pay out reserves, sell assets, borrow reserves; balance sheet shrinks\\
Bank rate & outside spread for discounting banks; internal drain absorbed by expanding the balance sheet\\
External drain & raise Bank rate (world rate of interest) or suspend specie payments\\
Without gold & both outside spreads vanish; CIP $F/S=(1+R^\ast)/(1+R)$ is the bridge to the modern system\\
\end{tabularx}
\end{keyformulas}

\section*{Exercises}

\begin{exercise}[A bill from Buenos Aires]
An Argentine exporter sells wheat to a German importer against a 90-day bill for \pounds1000, discounted in London at 4\% a year. Write all the T-accounts on the day
of discount and 90 days later, including the two FX dealers. How much sterling does the exporter receive?
\end{exercise}

\begin{exercise}[Gold points]
The dollar is 0.04838 ounces of gold and the pound 0.2354 ounces; shipping gold costs 0.6\% of its value. Compute the mint parity in pounds per dollar and the gold
points. Draw the FX dealer's diagram when the US is a deficit country.
\end{exercise}

\begin{exercise}[Three ways to defend]
A central bank holds 50 of gold, 100 of T-bills, and has 150 of currency outstanding. Dealers bring 60 of currency for gold. Show the balance sheet after each of
the three defenses, and say which ones raise domestic interest rates.
\end{exercise}

\begin{exercise}[Internal and external drain]
Explain why the Bank of England can stop an internal drain by lending freely at Bank rate but must raise Bank rate to stop an external drain. What does suspension
change for a country that is not at the top of the hierarchy?
\end{exercise}
